\documentclass[10pt,journal,letterpaper]{IEEEtran}
\usepackage[T1]{fontenc}
\usepackage[utf8]{inputenc}
\usepackage{amsmath,amssymb,booktabs,array,graphicx}
\usepackage{cite}
\usepackage[hidelinks]{hyperref}
\newcolumntype{P}[1]{>{\raggedright\arraybackslash}p{#1}}
\newcommand{\RevisionCases}{137}
\newcommand{\RevisionDisagreements}{0}

\begin{document}
\title{Engineering Semantics for Grounded Industrial Agents: An Equal-Information Evaluation}
\author{Anand~Todkar, Jitendra~Solanki, Mrinmoy~Sarkar}
\maketitle
\begin{abstract}
Industrial agents must associate requests with the correct equipment, interpret
unit-bearing observations and propose commands consistent with current machine
state. We investigate how explicit engineering facts support these tasks, while
separating their information content from representation and validation engine.
An agent-usability framework connects engineering capabilities to operational
failure modes. On a synthetic reference plant, an equally informed procedural
validator and a declarative SHACL validator agree on all \RevisionCases{}
proposed actions, each detecting every injected fault without rejecting legal
actions. Thus this corpus supports the value of the supplied engineering facts,
not an accuracy advantage for SHACL over competent procedural code. Additional
controls cross retrieval strategies with context conditions, remove individual
capabilities, probe missing evidence and check serialisation fidelity. A physical
S7-1500/industrial-PC testbed protocol is specified but has not yet been executed.
Transfer effort, operational latency and physical outcomes remain open evidence
requirements rather than claimed results.
\end{abstract}
\begin{IEEEkeywords}
Industrial agents, semantic interoperability, OPC UA, action validation,
engineering knowledge graphs, programmable logic controllers.
\end{IEEEkeywords}

\noindent\textit{Working revision, not submission-ready. Physical trials,
independent annotation, AUF reliability, specific competitor verification and
the current journal page policy remain pending.}

\section{Introduction}
A request to change a machine setting contains several distinct obligations:
identify the machine and writable signal, interpret the requested unit, obtain
current evidence, and respect operating-state and process guards. These are
supervisory decisions above the controller's deterministic loop. An agent's
plausible response is not itself evidence that a proposed command is admissible.

Industrial standards already describe much of this information. The problem
examined here is how to expose and check those facts at agent runtime, not whether
RDF is intrinsically safer than an information model or procedural implementation.
OPC UA and AAS can carry richer profiles than a flat typed catalogue; our
experimental conditions are selected projections, not upper bounds on the
standards' capabilities \cite{iec62541,idta2023aas}.

We pursue three contributions: (1) a capability framework linking engineering
evidence to agent failures; (2) an operational integration pattern for retrieval,
typed tools and pre-dispatch validation; and (3) controlled evaluation separating
missing information, retrieval limitations and implementation choice. The third
contribution currently has synthetic evidence only. Industrial transfer and
physical execution must be evaluated before deployment claims are made.

\section{Positioning and Scope}
RAMI 4.0 is an industrial reference architecture, whereas the proposed view
organises the evidence needed for individual agent decisions. They are
complementary: information and functional concerns provide engineering inputs;
communication concerns deliver observations and commands; integration and asset
concerns anchor physical identity; business concerns supply authorisation and
policy. This is a conceptual crosswalk, not a normative mapping or demonstrated
improvement over RAMI. The exact RAMI edition and source must be verified before
submission.

The five survey layers classify identification, information models, formal
semantics, domain vocabularies and cross-organisational policy. They are not a
strict dependency hierarchy. Context conditions C0--C4 and validator tiers
G0--G4 are experimental treatments assembled from these sources, not the survey
layers themselves. Dataspace policy is not evaluated here.

Graph retrieval and knowledge-graph augmentation provide relevant prior art
\cite{edge2024graphrag,pan2024unifying}. SHACL is an existing validation language,
not our invention \cite{knublauch2017shacl}. The contribution cannot be the
generic combination of an LLM and a graph. Specific neurosymbolic industrial
systems, including the suggested CausalTrace comparator, require source and
artefact verification; no capability absence or superiority is asserted before
that review. The quoted 94\% MAP@3 result is not used without a verified dataset,
ranked-cause task and common evaluation protocol.

\section{Capability Framework and Operational Method}
The agent-usability framework (AUF) separates identity (U1), types and access
(U2), quantities and conversion (U3), relations (U4), behaviour (U5), constraints
(U6), provenance and time (U7), query pushdown (U8), context economy (U9), and
ecosystem/deployability (U10). These are complementary capabilities; cumulative
ablation does not establish mathematical orthogonality.

The practical distinction is between facts a source can express, facts a
deployment actually publishes and facts available at decision time. A unit
conversion requires both source and target definitions or a documented external
conversion service. State legality requires a current state and transition model.
An interlock requires its guard, dependencies and observations. Absent evidence
is distinct from evidence of an illegal command.

The proposed runtime resolves the target, fetches the relevant evidence, generates
a typed tool proposal and validates it. An incomplete evidence set produces
abstention or a request for refresh. A proposal permitted at planning time is
checked again at dispatch; the PLC retains independent, current-state checks and
final actuator authority. A finite retry budget, operator approval and an audit
trail are required for repair. No LLM participates in a hard-real-time loop.

AUF scores remain expert assessments until independently checked. The planned
assessment uses two scorers, specification-clause evidence, ordinal agreement
and adjudication, then tests whether AUF-guided capability selection improves a
held-out installation within a fixed integration budget. Those results are not
available in this revision.

\section{Controlled Evaluation}
The existing generated reference plant contains 17 assets and 73 signals. Its
35 tasks and 137 proposed actions are retained for comparability, not presented
as independently collected industrial workloads. All revision outputs are
isolated from the original manuscript and archived results.

\subsection{R1: Equal-information validation}
Both implementations receive the same engineering ranges, unit calculus,
asset containment, PackML state/commands, interlocks and observations. The
procedural implementation reads those declarative reference data without calling
SHACL. Predictions do not inspect gold labels. Raw constraint-family sets are
compared; diagnosis maps an injected scale error to a range violation identically
for both implementations, rather than adding labels to predictions from gold.
Detection, false positives and diagnosis are reported separately. Binomial Wilson
intervals are archived but are descriptive for this designed corpus, not evidence
of installation-wide incident prevalence.

\subsection{R2: Retrieval crossed with context}
Lexical retrieval is evaluated under all conditions. Membership-scoped retrieval
is added where published asset membership is available; both scoped strategies
use a fixed manually supplied entity resolver. Graph traversal is restricted to
C3/C4. Each supported strategy is tested at 1,500, 6,000 and 24,000 surrogate
tokens with a fixed condition renderer. An unbudgeted all-elements context is an
oracle diagnostic, not a deployable retriever. No required-fact set is used for
ranking. This matrix separates some retrieval effects but does not equalise all
notation familiarity or entity-linking assumptions.

\subsection{R3--R5: Ablation, degradation and fidelity}
R3 removes units, flow topology, behaviour, constraints or provenance separately.
The context ordering remains fixed to full C4 as a diagnostic intervention;
it is not evidence that an ablated deployment could build that ordering. The
validator separately removes access to relevant checking capabilities and counts
abstentions, which must not be counted as detections. Context flow topology and
validator containment are different interventions and cannot be pooled as one
effect size.

R4 probes missing state, unknown/missing unit definitions, malformed timestamps,
missing observations and a synthetic planning-to-dispatch state change. R5
serialises one RDF graph and parses each syntax back, qualifying an
equal-information comparison only if graph isomorphism is preserved.

Evidence sufficiency is the fraction of tasks whose declared required facts are
present. It is not an unconditional accuracy bound: alternative derivations,
external knowledge and inferred conventions may change answerability. Conversion
facts are conservatively counted only when both unit definitions occur in the
context. Token counts are a deterministic surrogate, not model billing tokens.

\section{Results and Interpretation}
\begin{table}[t]
\centering\footnotesize
\caption{R1: equal-information validation on the synthetic reference corpus. Diagnosis maps scale faults to range violations for both implementations.}
\label{tab:revision-equal}
\begin{tabular}{@{}llll@{}}
\toprule
Validator & Recall & FPR & Diagnosis \\
\midrule
Procedural & 1.000 & 0.000 & 1.000 \\
SHACL & 1.000 & 0.000 & 1.000 \\
\bottomrule
\end{tabular}
\end{table}

R1 produces \RevisionDisagreements{} raw-family disagreements. Both implementations
detect every injected fault and accept every legal reference action
(Table~\ref{tab:revision-equal}). This is a negative result for an
engine-superiority claim and positive evidence that the supplied facts can be
checked independently. Standardisation, reuse and maintenance advantages need
separate effort measurements.

\begin{table}[t]
\centering\footnotesize
\caption{R2: crossed retrieval at 6,000 surrogate tokens. Oracle-all is unbudgeted and diagnostic only. Membership uses a fixed shared entity resolver.}
\label{tab:revision-retrieval}
\begin{tabular}{@{}llll@{}}
\toprule
Context & Retrieval & Suff. & Coverage \\
\midrule
C2 & lexical & 0.229 & 0.549 \\
C2 & oracle\_all & 0.314 & 0.649 \\
C2 & membership & 0.229 & 0.485 \\
C4 & lexical & 0.714 & 0.846 \\
C4 & oracle\_all & 0.971 & 0.995 \\
C4 & membership & 0.457 & 0.688 \\
C4 & graph & 0.800 & 0.929 \\
\bottomrule
\end{tabular}
\end{table}

Table~\ref{tab:revision-retrieval} distinguishes restricted contexts from
unbudgeted evidence. Under the audited conversion definition, a task may lack
an explicit target-unit definition even when the original bookkeeping marked
conversion as available. Differences from archived results are therefore not
silently treated as replication failures or attributed solely to routing.

\begin{table}[t]
\centering\footnotesize
\caption{R3: capability removal. Context sufficiency uses a fixed full-C4 ordering; abstentions are not detections. Context topology removes flow only; validator topology removes scope-containment checks.}
\label{tab:revision-ablation}
\begin{tabular}{@{}llll@{}}
\toprule
Removed & Suff. & Recall & Abstain \\
\midrule
none & 0.800 & 1.000 & 0 \\
behaviour & 0.629 & 0.893 & 37 \\
constraints & 0.686 & 0.964 & 20 \\
provenance & 0.743 & 0.845 & 40 \\
topology & 0.629 & 0.940 & 5 \\
units & 0.800 & 0.607 & 60 \\
\bottomrule
\end{tabular}
\end{table}

Removal effects depend on the workload and budget
(Table~\ref{tab:revision-ablation}). Smaller contexts can offset lost capability
for some tasks; lack of a net sufficiency decrease does not prove a capability
unnecessary. Additional per-category and per-task outputs are archived. The
corpus was designed around these failure families, so general necessity requires
independently authored tasks and external workloads.

\begin{table}[t]
\centering\footnotesize
\caption{R4: synthetic missing-evidence and dispatch-time state probes. These are not PLC trials or physical safety measurements.}
\label{tab:revision-degradation}
\begin{tabular}{@{}ll@{}}
\toprule
Probe & Disposition \\
\midrule
before\_state\_change & accept \\
after\_state\_change & reject \\
missing\_state & abstain \\
missing\_unit & reject \\
unknown\_unit & reject \\
malformed\_time & reject \\
missing\_observation & abstain \\
\bottomrule
\end{tabular}
\end{table}

The synthetic state probe permits a volume change in Idle and rejects it after
a change to Execute. It demonstrates the test harness's revalidation rule, not
that a real PLC/agent pipeline prevents a timing race. Missing evidence remains
visible in the verdict even when a separately proven fault also requires rejection.

\begin{table}[t]
\centering\footnotesize
\caption{R5: identical RDF graph round-tripped through three syntaxes. Only isomorphic round trips qualify as equal-information controls. Counts use the surrogate tokenizer.}
\label{tab:revision-serialisation}
\begin{tabular}{@{}llll@{}}
\toprule
Syntax & Triples & Tokens & Isomorphic \\
\midrule
turtle & 2465 & 45803 & no \\
nt & 2465 & 179200 & yes \\
json-ld & 2465 & 145212 & yes \\
\bottomrule
\end{tabular}
\end{table}

The installed RDFLib Turtle serializer rounds some double literals, causing a
strict graph mismatch despite identical triple counts. Such output is excluded
from equal-information claims; N-Triples and JSON-LD qualify in this run.
Preserving datatype and numeric precision is a deployment requirement, not
merely a token-economy detail.

\section{Physical Testbed Protocol: Pending}
We propose an S7-1500 controlling two guarded 24~V handling stations, optionally
coupled to a low-head water-filling module with physical level/flow sensors. An
IPC427E hosts the agent and semantic/validation services; an IPC227G can provide
an edge gateway or separate disturbance/replay service. Actual CPU order number,
firmware, OPC UA support/licence, I/O and IPC resources must be recorded before
commissioning. A PLC with simulated I/O is hardware-in-the-loop evidence, not
physical process validation.

Three workflows exercise unit-correct reporting, admissible setpoint/state
proposals and incident diagnosis. Baselines receive identical available
engineering evidence and tool permissions. Tasks and root-cause labels are
reviewed by an independent engineer; assets or operating episodes, not question
paraphrases, define held-out splits. Trials begin in read-only shadow mode.
Invalid proposals are never sent to actuators. Approved legal proposals are
bounded by PLC-side guards, operator supervision and an independently assessed
stop function. A standard CPU is not presumed safety-rated.

Primary outcomes are grounded-answer correctness, invalid proposals admitted,
legal proposals rejected and complete legal task execution. Record observation
age, validation/dispatch latency, controller acknowledgement and actual process
response separately. Real-time claims require an agreed supervisory deadline
and measured tail latency; the synthetic 15-minute cutoff is not a physical
freshness setting. A held-out station/configuration measures mapping edits,
custom code changes, regression failures and engineer-hours. Full protocol and
recording requirements are in the revision artefact; no trial results are claimed.

\section{Limitations and Deployment Implications}
The benchmark has one generated plant, authored tasks and injected faults.
Equal-information procedural validation shares the reference rules and does not
provide independent truth. Repeated model runs must be analysed within episodes,
not treated as independent incidents. Corrupted constraints, incorrect mappings,
incomplete interlock inventories, unseen fault combinations and future-dated
evidence remain additional threats. The current harness is not a live dispatcher.

Physical evaluation must distinguish agent proposal quality from gateway rejection
and PLC protection. Zero observed invalid actions cannot establish a safety
integrity level. Dataspace permission is a separate untested extension: ODRL
expresses policies, while trusted enforcement and downstream-use control require
additional mechanisms; it is not automatically interchangeable with SHACL.

\section{Conclusion}
The controlled corpus supports engineering-evidence checks, not an inherent
accuracy advantage of semantic graphs over an equally informed rule program.
The proposed framework is useful only insofar as it helps expose, select and
maintain the facts industrial decisions require. Physical performance, independent
capability assessment and transfer effort are the next evidence gates.

\bibliographystyle{IEEEtran}
\bibliography{../paper/refs}
\end{document}